\chapter{Work, energy and power}\label{ch:energy}

Force is local; energy is the accounting trick that lets you skip the
details. $W=\int\mathbf F\cdot d\mathbf r$ accumulates what forces deposit,
$K=\tfrac12mv^{2}$ is the currency it arrives in, and potential energy is
the bank where conservative forces store it. The chapter's payoff is the
right to answer ``how fast'' without ever solving for the trajectory.

% from-zero:begin
\begin{rulebox}[Before you start]
Chapters~1--3. \emph{Recap:} work is force times distance along the motion; energy is the bank account that work pays into.
\end{rulebox}
\begin{wins}[After this chapter you can]
\begin{itemize}
\item compute work for constant and varying forces;
\item use the work--energy theorem and conservation of energy, and know when conservation fails;
\item read motion from an energy diagram;
\item do an energy audit of a real system.
\end{itemize}
\end{wins}
\begin{principle}[How to think about this]
Energy is bookkeeping for motion: it changes form but is never lost, and when it seems to disappear it has gone into heat or sound. Often the shortest route to a speed or height is to compare energy at two moments and skip the forces entirely.
\end{principle}
% from-zero:end

\section{Work: force collected along a path}\label{sec:en-work}

\begin{equation}
  W=\int\mathbf F\cdot d\mathbf r.\label{eq:work}
\end{equation}
Work is a \emph{path integral}: only the force component along the motion
counts. The centripetal force of \S1.5 does \emph{no} work --- always
perpendicular, always at right angles to $d\mathbf r$, which is why
uniform circular motion keeps its speed. Figure~\ref{fig:work} is the
one-dimensional version: area under $F(x)$.

\begin{figure}[htbp]\centering
\begin{tikzpicture}
\begin{axis}[width=8.5cm,height=5cm,xlabel={$x$ (m)},ylabel={$F$ (N)},
xmin=0,xmax=10,ymin=0,ymax=6,xtick={0,2,4,6,8,10},font=\small]
\addplot[fill=accent!15,draw=none,domain=0:8] {4 - 0.3*x} \closedcycle;
\addplot[very thick,accent,domain=0:8] {4 - 0.3*x};
\node[font=\scriptsize] at (axis cs:3.5,1.2) {$W=\displaystyle\int F\,dx$};
\end{axis}
\end{tikzpicture}

\caption{Work as area under the $F(x)$ curve. Notice variable forces need
the integral; the constant-force $W=Fd\cos\theta$ is only its rectangle.}
\label{fig:work}
\end{figure}

\section{The work--kinetic-energy theorem}\label{sec:en-wke}

Start from $\mathbf F=m\,d\mathbf v/dt$ and dot with $d\mathbf r$:
\[
  \mathbf F\cdot d\mathbf r
   =m\frac{d\mathbf v}{dt}\cdot d\mathbf r
   =m\,\mathbf v\cdot d\mathbf v
   =d\!\left(\tfrac12mv^{2}\right),
\]
so
\begin{equation}
  W_{\rm net}=\Delta K.\label{eq:wke}
\end{equation}
Every joule of net work reappears exactly as kinetic energy --- the
theorem is a tautology dressed as a tool, and it works for any path,
including ones too ugly to integrate force over directly.

\section{Potential energy: where conservative forces bank work}\label{sec:en-pe}

A force is \emph{conservative} when the work it does depends only on the
endpoints. Then you can declare a stored quantity:
$\Delta U=-W_{\rm by\ force}$. The two workhorses:
\[
  U_g=mgy,\qquad U_s=\tfrac12kx^{2}.
\]
Friction fails the test --- the longer the path, the more it takes --- so
there is no ``potential energy of friction''; its work leaves the
mechanical books entirely (section~\ref{sec:en-noncons}).

\begin{figure}[htbp]\centering
\begin{tikzpicture}
\begin{axis}[width=8.5cm,height=5.5cm,xlabel={$x$},ylabel={$U(x)$},
xmin=-0.5,xmax=3,ymin=0,ymax=4.5,xtick={0,1,2,3},font=\small]
\addplot[very thick,accent,domain=0:2.6] {0.5*x^2};
\addplot[dashed,amber] coordinates {(0,2) (2.6,2)};
\addplot[mark=*] coordinates {(2,2)} node[above,font=\scriptsize]{turning};
\draw[->,thick] (axis cs:2,1.55) -- (axis cs:1.55,0.9);
\draw[->,thick] (axis cs:2,1.55) -- (axis cs:2.45,0.9);
\node[font=\scriptsize,amber] at (axis cs:1.2,2.25) {$E$ fixed};
\node[font=\scriptsize] at (axis cs:2.35,0.55) {$K>0$ region};
\end{axis}
\end{tikzpicture}

\caption{Reading a $U(x)$ diagram. With $E$ fixed, motion lives where
$K=E-U>0$; the turning point is where the parabola meets the line. The
restoring force is $-dU/dx$, the slope reversed.}
\label{fig:ux}
\end{figure}

% from-zero:naive-carry:begin
\begin{example}[First attempt: multiplying force by distance without checking direction]\label{ex:naive-carry}
You carry a $20\,\mathrm{kg}$ box $10\,\mathrm{m}$ across a level floor at constant speed. A quick estimate of the work you
did: $W=mgd=20\times9.8\times10=1960\,\mathrm{J}$.
\par\smallskip
\textbf{What goes wrong.} Work is $W=\mathbf F\cdot\mathbf d=Fd\cos\theta$. The force you apply to hold the box up is
vertical ($mg=196\,\mathrm{N}$, upward) and the displacement is horizontal, so $\theta=90^\circ$ and
$\cos\theta=0$: you do \emph{no} work on the box. (Your muscles still burn energy, which is a different matter
entirely.) The $1960\,\mathrm{J}$ answer would be right if you lifted the box $10\,\mathrm{m}$ straight up.
\end{example}
% from-zero:naive-carry:end

\section{Conservation of mechanical energy --- and its failure modes}\label{sec:en-cons}

When only conservative forces work,
\begin{equation}
  K_i+U_i=K_f+U_f.\label{eq:emech}
\end{equation}
Figure~\ref{fig:bars} is the account statement: the bars trade height,
never total. The failure modes are equally instructive: friction converts,
deformation absorbs, and anything with a motor adds from outside --- all
of them break equation~\eqref{eq:emech}, all handled by
section~\ref{sec:en-noncons}'s audit.

\begin{figure}[htbp]\centering
\begin{tikzpicture}[scale=0.85,font=\scriptsize]
% bar charts: launch, top, landing
\foreach \x/\lab/\K/\U in {0/launch/4/0, 3.4/apex/0/4, 6.8/landing/4/0}{
  \begin{scope}[xshift=\x cm]
    \draw[thick] (0,0) rectangle (2,2.6);
    \fill[accent!70] (0,0) rectangle (0.9,\K*0.55);
    \fill[amber!60] (1.1,0) rectangle (2,\U*0.55);
    \node at (1,-0.4) {\lab};
    \node[rotate=90] at (0.45,\K*0.28) {K};
    \node[rotate=90] at (1.55,\U*0.28) {U};
  \end{scope}}
\node at (3.4,3.3) {$E=K+U$ constant: the bars trade, never grow};
\end{tikzpicture}

\caption{Energy bars along a trajectory: $K$ and $U$ trade, the sum
doesn't move.}
\label{fig:bars}
\end{figure}

\begin{example}[Minimum release height for the loop-the-loop]\label{ex:loop}
A cart starts from rest at height $h$ on a frictionless track containing a
loop of radius $r$ (figure~\ref{fig:loop}). Minimum $h$ to make it
around?
\par\smallskip
\textbf{Strategy.} Two chapters meet: energy sets speeds
(equation~\eqref{eq:emech}), circular dynamics (ch2) sets the requirement
--- at the top, $N\ge0$ means $v^{2}\ge gr$.
\begin{solution}
Top of loop is at height $2r$; borderline case has $v_{\rm top}^{2}=gr$.
Energy: $mgh=mg(2r)+\tfrac12m\,gr$, so
\[
  h=\frac{5}{2}r.
\]
Real tracks need more --- friction and the cart's own rotation (ch5) both
take a share. \checkmark
\end{solution}
\end{example}

\begin{figure}[htbp]\centering
\begin{tikzpicture}[scale=1.0]
\draw[thick] (-3.5,0) -- (0,0);
\draw[thick] (0,0) arc (270:180+180:1.6);
\draw[thick] (1.6,0) arc (270:360-270:1.6);
% circle at center (1.6,1.6) radius 1.6
\draw[thick] (1.6,1.6) circle (1.6);
\filldraw[fill=gray!20] (-1.2,0.9) rectangle (-0.7,1.4);
\draw[thick,dashed] (-1.2,0.9) -- (-3.0,0);
\draw[->,thick,accent] (-3.3,-0.3) -- (-0.8,-0.3) node[right,font=\scriptsize]{release at $h$};
\filldraw[fill=amber!50] (1.45,3.15) rectangle (1.75,3.45);
\node[font=\scriptsize] at (2.6,3.3) {top: needs $v^{2}\ge gr$};
\draw[dashed] (-0.5,3.3) -- (3.4,3.3) node[right,font=\scriptsize]{$2r$};
\draw[dashed] (-0.5,1.6) -- (1.6,1.6);
\node[font=\scriptsize] at (-0.2,1.35) {$r$};
\end{tikzpicture}

\caption{The loop-the-loop energy budget. The $5r/2$ answer is the
classic meeting point of ch1's $a_c=v^{2}/r$ and this chapter's
bookkeeping.}
\label{fig:loop}
\end{figure}

\begin{example}[Spring-launched block up a rough incline]\label{ex:springincl}
A spring ($k=200$\,N/m) compressed $0.3$\,m launches a $0.5$\,kg block up
a $\theta=30^{\circ}$ incline with $\mu_k=0.2$. How far along the slope
does it go?
\par\smallskip
\textbf{Strategy.} Energy audit with friction present: spring's deposit
minus friction's withdrawal equals the height gained.
\begin{solution}
$U_s=\tfrac12(200)(0.09)=9$\,J available. Along distance $d$: gravity
takes $mgd\sin30^{\circ}=2.45d$\,J, friction takes
$\mu_k mg\cos30^{\circ}\cdot d=0.849d$\,J. So
$d=9/(2.45+0.849)=2.73$\,m --- vs $3.67$\,m frictionless; friction's
share is a quarter of the budget. \checkmark
\end{solution}
\end{example}

\begin{example}[Pendulum speed at the bottom]\label{ex:pendulum}
A $0.5$\,kg mass on a $1$\,m string is released from horizontal. Speed at
the bottom?
\par\smallskip
\textbf{Strategy.} Pure energy --- the tension does no work (always
perpendicular to the arc), so mechanical energy conserves.
\begin{solution}
$mgL=\tfrac12mv^{2} \Rightarrow v=\sqrt{2gL}=\sqrt{19.6}=4.43$\,m/s ---
independent of mass and of the path taken to get there. The same
$\sqrt{2gh}$ a frictionless slide gives: energy doesn't care about
trajectories. \checkmark
\end{solution}
\end{example}

\begin{exercise}[easy]
Find the pendulum's speed when the string passes $45^{\circ}$ from
vertical. (AP-C)
\end{exercise}

\section{Energy diagrams: motion read off a picture}\label{sec:en-diagrams}

Given $U(x)$ (figure~\ref{fig:ux}): total energy $E$ draws a horizontal
line; kinetic energy is the gap below it; turning points are where the
line crosses the curve; equilibria sit at the curve's extrema (stable at
minima, unstable at maxima). A particle in a quartic double well can
\emph{never} leave one side if its line dips below the saddle --- that is
the whole qualitative solution, no ODE required.

\section{Power: the rate of doing work}\label{sec:en-power}

\begin{equation}
  P=\frac{dW}{dt}=\mathbf F\cdot\mathbf v.\label{eq:power}
\end{equation}
A $1000$\,kg car climbing a $6\%$ grade at $20$\,m/s burns
$P=mg\sin\theta\cdot v\approx11.8$\,kW on gravity alone --- before drag's
share, which figure~\ref{fig:power} adds.

\begin{figure}[htbp]\centering
\begin{tikzpicture}
\begin{axis}[width=8cm,height=5cm,xlabel={$v$ (m/s)},ylabel={$P=Fv$ (kW)},
xmin=0,xmax=30,ymin=0,ymax=40,font=\small]
\addplot[very thick,accent,domain=0:30] {0.04*x + 0.01*x^2};
\node[font=\scriptsize,align=center] at (axis cs:14,26) {rolling $+$ drag:\\$P=(F_{\rm roll}+cv)v$};
\addplot[dashed,amber] coordinates {(20,0) (20,12)};
\node[font=\scriptsize,amber,align=center] at (axis cs:21,6) {$72$ km/h:\\$\sim12$ kW};
\end{axis}
\end{tikzpicture}

\caption{$P=\mathbf F\cdot\mathbf v$: climbing power is linear in $v$,
drag power quadratic-ish. Notice how quickly the sum leaves the
city-driving budget.}
\label{fig:power}
\end{figure}

\begin{example}[Stair-climb power measurement]\label{ex:stairs}
You ($65$\,kg) run up a $4$\,m stair flight in $5$\,s. Estimate your
mechanical power output.
\par\smallskip
\textbf{Strategy.} $P=\Delta U/t$ --- the simplest honest measurement in
the book.
\begin{solution}
$\Delta U=mgh=65\times9.8\times4=2548$\,J, so
$P\approx510$\,W $\approx0.7$\,hp. A cyclist sustains $\sim250$\,W; elite
sprinters peak near $2$\,kW for seconds. \checkmark
\end{solution}
\end{example}

\begin{exercise}[medium]
A pump lifts $200$\,L/min of water $10$\,m. Minimum motor power? (F=ma)
\end{exercise}

\section{Elastic collisions: the deferred full treatment}\label{sec:en-elastic}

Chapter 3 owned momentum; energy completes the elastic case. The
$v_{\rm rel}$-flip shortcut: in an elastic collision the relative velocity
\emph{reverses} --- $v_{2f}-v_{1f}=-(v_{2i}-v_{1i})$ --- which with
momentum alone solves everything.

\begin{example}[Elastic collision, three ways]\label{ex:elastic}
$m_1=2$\,kg at $6$\,m/s hits stationary $m_2=2$\,kg. Solve by (a)
lab-frame algebra, (b) the CM frame, (c) the flip rule.
\par\smallskip
\textbf{Strategy.} Three routes to the same answer --- each teaches a
different skill.
\begin{solution}
(a) $v_1'+v_2'=6$ (momentum) and $v_1'^{2}+v_2'^{2}=36$ (energy):
substituting gives $v_1'v_2'=0$, so one stops, one carries $6$\,m/s ---
velocities swap.
(b) CM moves at $3$\,m/s; in its frame the blocks approach at $\pm3$ and
must leave at $\mp3$ (reversal is the only $K$-preserving option).
Transforming back: $0$ and $6$.
(c) Flip: $v_2'-v_1'=6$ plus $v_1'+v_2'=6$ gives $v_1'=0$, $v_2'=6$.
All three agree; the equal-mass swap is figure~\ref{fig:pk}. \checkmark
\end{solution}
\end{example}

\begin{figure}[htbp]\centering
\begin{tikzpicture}[scale=0.9,font=\scriptsize]
\begin{scope}
\node at (1.5,2.2) {before};
\draw[->,thick,accent] (0,1) -- (1,1) node[right]{$v$};
\draw[fill=gray!20] (-0.5,0.6) rectangle (0,1.4);
\draw[fill=gray!20] (2.2,0.6) rectangle (2.7,1.4);
\node at (-0.25,0.4) {$m$}; \node at (2.45,0.4) {$m$};
\node at (1.5,-0.3) {$p=mv$,\quad $K=\tfrac12mv^{2}$};
\end{scope}
\begin{scope}[xshift=4.6cm]
\node at (1.5,2.2) {after ($e=1$)};
\draw[fill=gray!20] (0,0.6) rectangle (0.5,1.4);
\draw[->,thick,accent] (2.7,1) -- (3.7,1) node[right]{$v$};
\draw[fill=gray!20] (2.2,0.6) rectangle (2.7,1.4);
\node at (0.25,0.4) {$m$}; \node at (2.45,0.4) {$m$};
\node at (1.7,-0.3) {$p=mv$,\quad $K=\tfrac12mv^{2}$};
\end{scope}
\node[align=left] at (4.6,-1.6) {equal masses: velocities swap ---\\both $p$ and $K$ survive};
\end{tikzpicture}

\caption{Bookkeeping for the equal-mass elastic hit: both $p$ and $K$
survive --- a rare event, reserved for $e=1$.}
\label{fig:pk}
\end{figure}

\begin{example}[General elastic formula]\label{ex:elasticgen}
Apply the flip rule to unequal masses: $m_1=5$\,kg at $8$\,m/s hits
$m_2=1$\,kg at rest.
\par\smallskip
\textbf{Strategy.} Two equations: momentum and the reversed relative
velocity.
\begin{solution}
$5v_1'+v_2'=40$ and $v_2'-v_1'=8$:
subtracting, $6v_1'=32 \Rightarrow v_1'=5.33$\,m/s,
$v_2'=13.33$\,m/s. Check $K$: $\tfrac12(5)(28.4)+\tfrac12(1)(177.7)=
71+88.9=160$\,J $=\tfrac12(5)(64)$. \checkmark
\end{solution}
\end{example}

\section{Non-conservative work: the energy audit}\label{sec:en-noncons}

With friction present, equation~\eqref{eq:emech} becomes the audit
\[
  K_i+U_i+W_{\rm nc}=K_f+U_f,
  \qquad W_{\rm nc}=-f_kd<0 \ \text{(usually)}.
\]
The missing energy isn't gone --- it's heat, sound, and deformation. The
conveyor belt is the classic trap: dragging a block up to belt speed
heats the belt-block interface by exactly the kinetic energy the block
gains ($Q=K$) --- half the motor's work becomes motion, half becomes
warmth.

\begin{lstlisting}[language=Python,caption={\texttt{code/ch04\_energy\_check.py} --- plain Euler vs symplectic Euler on a block-spring.},label={lst:edrift}]
if scheme == "euler":
    x_new = x + v * dt
    v = v - (k/m) * x * dt        # BOTH updates use OLD values
    x = x_new
else:                            # symplectic
    v += -(k/m) * x * dt          # update v first
    x += v * dt                   # then x with the NEW v
\end{lstlisting}

Output of a real run ($k=4$, $m=1$, released from $x=1$):
\begin{lstlisting}[numbers=none]
E0 = 2.0000 J
scheme           dt  periods       E_end        drift
euler           0.1       10     446168.12 22308306.0%
euler          0.05       10       1034.79    51639.6%
symplectic      0.1       10          1.97       -1.5%
symplectic     0.05       10          2.00        0.1%
euler 100 periods, dt=0.1: E=446168.120 J (x 223084.1 of original -- secular growth)
sympl 100 periods, dt=0.1: E=1.971 J (bounded oscillation, no secular drift)
\end{lstlisting}

The checkpoint's prediction is now a number: plain Euler's error compounds
multiplicatively (the orbit spirals \emph{out}), while symplectic Euler ---
update $v$ before $x$ --- conserves a nearby shadow Hamiltonian and merely
oscillates around the true energy. Every orbit simulation after this
chapter uses the symplectic variant.

\begin{lstlisting}[language=Python,caption={\texttt{code/ch04\_energy\_budget.py} --- one car, three energy lines: KE, a climb's potential, drag's running cost.},label={lst:budget}]
m, v, g = 1200.0, 27.0, 9.81        # kg, m/s (~60 mph)
Cd, A, rho = 0.30, 2.2, 1.2
KE = 0.5 * m * v**2
PE_100m = m * g * 100
F_drag = 0.5 * rho * Cd * A * v**2
P_drag = F_drag * v
print(f"KE at 60 mph            : {KE/1e3:8.1f} kJ")
print(f"PE climbing 100 m       : {PE_100m/1e3:8.1f} kJ")
print(f"drag force              : {F_drag:8.0f} N")
print(f"drag power at 60 mph    : {P_drag/1e3:8.1f} kW")
print(f"1 min of drag work      : {P_drag*60/1e3:8.1f} kJ = {P_drag*60/KE:.2f} x KE")
print(f"double the speed: drag power x{2**3}")
\end{lstlisting}

Output of a real run:
\begin{lstlisting}[numbers=none]
KE at 60 mph            :    437.4 kJ
PE climbing 100 m       :   1177.2 kJ
drag force              :      289 N
drag power at 60 mph    :      7.8 kW
1 min of drag work      :    467.7 kJ = 1.07 x KE
double the speed: drag power x8
\end{lstlisting}

The audit surprises: a $100$\,m climb costs nearly $3\times$ the
car's entire highway kinetic energy --- hills are the expensive part
of the map --- while one minute of cruise drag burns the KE outright.
The last line is the speed trap: doubling $v$ makes drag \emph{power}
eight-fold, which is exactly why fuel economy falls off a cliff above
$90$\,km/h.

\section{In engineering practice: the energy audit}\label{sec:en-eng}

Regenerative braking is this chapter's audit run as accounting: an
electric car descending a hill deposits into the battery at $70$--$85\%$
efficiency --- not $100\%$, because the audit of
section~\ref{sec:en-noncons} still applies (aerodynamic drag, inverter
losses, battery chemistry each take a line). Roller-coaster designers run
example~\ref{ex:loop} backward: given a required top speed, the lift hill
is $h=5r/2$ plus an empirically tuned friction allowance --- which is why
first hills tower over every loop.

\section{Common mistakes}\label{sec:en-mistakes}

\begin{tripwire}[Sign of work vs $\Delta U$]
$W_{\rm by}=+\Delta K$; $\Delta U=-W_{\rm by\ conservative}$. Mixing
``work done on'' with ``work done by'' flips every sign in the audit.
\end{tripwire}
\begin{tripwire}[Mixing PE zero points]
$U=mgh$ measures against \emph{your} chosen floor; two different floors in
one problem corrupt the difference.
\end{tripwire}
\begin{tripwire}[Conserving mechanical energy through friction]
Friction withdraws before gravity deposits --- example~\ref{ex:springincl}
lost a quarter of its budget to it.
\end{tripwire}
\begin{tripwire}[$P=Fv$ with the wrong $v$]
The velocity in $P=\mathbf F\cdot\mathbf v$ is the \emph{point of
application's} --- for a moving belt problem it's the belt, not the block.
\end{tripwire}
\begin{tripwire}[``Contact forces do no work'' on deformables]
True for rigid bodies; false for cars, balls, and people --- deformation
absorbs real joules.
\end{tripwire}

\section{Summary}\label{sec:en-summary}

\begin{table}[ht]\centering
\begin{tabularx}{\linewidth}{@{}lY@{}}
\toprule
\tablehead\thead{Idea} & \thead{One-line version}\\
\midrule
work & $W=\int\mathbf F\cdot d\mathbf r$; only the along-motion part counts \\
WKE theorem & $W_{\rm net}=\Delta K$ --- derived, not assumed \\
potential energy & $\Delta U=-W_{\rm cons}$; $mgy$ and $\tfrac12kx^{2}$ \\
conservation & $K+U$ constant iff only conservative forces work \\
energy diagram & turning points where $E$ meets $U$; equilibria at $U'$ extrema \\
power & $P=\mathbf F\cdot\mathbf v$ \\
elastic collisions & relative velocity flips; with $p$ it solves everything \\
energy audit & $W_{\rm nc}$ bridges the books; friction $\Rightarrow$ $-f_kd$ \\
\bottomrule
\end{tabularx}
\caption{Energy cheat sheet.}
\label{tab:ensummary}
\end{table}

\begin{wins}[Definition of done]
You can now: integrate variable forces; derive the work--energy theorem;
read motion off $U(x)$; size a loop-the-loop hill; audit friction's share;
solve elastic collisions three ways; and explain why symplectic
integrators run the book's orbits.
\end{wins}

\chapternotes{%
The work--energy ordering follows MIT 8.01SC weeks 7--9 \cite{ocw801} and
AP Physics C unit 3 \cite{apced}. OpenStax College Physics ch.~6--7
\cite{openstax1} has parallel treatment; Morin ch.~5 \cite{morin} is the
source of the energy-diagram exercises. The integrator comparison is the
standard motivation for symplectic schemes; loop-the-loop numbers are the
classic result.}

\begin{rulebox}[{Watch instead}]
\begin{enumerate}
\item \emph{Work, Kinetic \& Potential Energy, Gravitation, Conservative
Forces} --- Walter Lewin, MIT 8.01x Lecture 11, 49:06,
\url{https://youtu.be/9gUdDM6LZGo} --- \S\ref{sec:en-work}--\S\ref{sec:en-cons}
at full length (accessed 2026-10-11).
\item \emph{Work-energy theorem} --- Khan Academy India, 12:34,
\url{https://youtu.be/nnMirzeTepQ} --- the §\ref{sec:en-wke} derivation
worked quickly (accessed 2026-10-11).
\end{enumerate}
\end{rulebox}

\section*{Further reading}
\addcontentsline{toc}{section}{Further reading}

\begin{itemize}
\item MIT OpenCourseWare, \emph{8.01SC Classical Mechanics}, weeks 7--9 ---
work, energy and collisions with recitation solutions. At
\url{https://ocw.mit.edu}, course ``8.01SC'' (verified 2026-10-11).
\item OpenStax, \emph{College Physics 2e}, ch.~6--7 ``Work, Energy, and
Energy Resources'' --- free text at \url{https://openstax.org} (verified
2026-10-11).
\item College Board, \emph{AP Physics C: Mechanics CED}, Unit 3 --- the
exam blueprint for this chapter. At
\url{https://apstudents.collegeboard.org} (verified 2026-10-11).
\item D.~Morin, \emph{Introduction to Classical Mechanics}, ch.~5 ---
energy methods at the harder level; ISBN 978-0521876223.
\item The Physics Classroom, \emph{Work, Energy, and Power} --- readable
remediation with roller-coaster walkthroughs, at
\url{https://www.physicsclassroom.com} (verified 2026-10-11).
\item Khan Academy, \emph{Work and energy} --- drills for this whole
chapter, at \url{https://www.khanacademy.org} under Science $\to$ Physics
(verified 2026-10-11).
\end{itemize}

\section*{Additional exercises}
\addcontentsline{toc}{section}{Additional exercises}

\begin{exercise}[easy]
Sign reasoning: classify the work of gravity, normal, and friction as
$+$, $-$, or $0$ for (a) a block sliding down an incline, (b) carried up
the same incline at constant speed, (c) a satellite in circular orbit.
(AP-C)
\end{exercise}
\begin{exercise}[easy]
A crane lifts a $500$\,kg load at $0.5$\,m/s steadily. Find the motor's
power. (AP-C)
\end{exercise}
\begin{exercise}[easy]
A spring ($k=400$\,N/m) stores how much energy at $0.1$\,m compression?
What mass could it lift $1$\,m with that energy? (AP-C)
\end{exercise}
\begin{exercise}[easy]
Read figure~\ref{fig:bars}: if the launch bar holds $8$\,J and the mass
is $0.5$\,kg, find the launch speed and apex height. (F=ma)
\end{exercise}
\begin{exercise}[easy]
A child slides down a $3$\,m frictionless slide. Find the exit speed,
then explain why the slide's shape doesn't matter. (AP-C)
\end{exercise}
\begin{exercise}[medium]
Loop minimum speed: recompute example~\ref{ex:loop} requiring
$v_{\rm top}^{2}=2gr$ (a safety margin) --- what $h$ now? (F=ma classic)
\end{exercise}
\begin{exercise}[medium]
Spring + rough incline: example~\ref{ex:springincl} again with
$k=300$\,N/m, $\mu_k=0.3$, $\theta=20^{\circ}$ --- find $d$, then the
block's speed back at the launch point after sliding down. (AP-C FR-style)
\end{exercise}
\begin{exercise}[medium]
Stair power: have a partner time you on a real flight; compute your watts
and compare against the $510$\,W estimate. (F=ma, data)
\end{exercise}
\begin{exercise}[medium]
A pendulum released from $60^{\circ}$ clears the bottom at what fraction
of the $90^{\circ}$-release speed? Work it in symbols first. (AP-C)
\end{exercise}
\begin{exercise}[medium]
Elastic collision: $m_1=3$\,kg at $10$\,m/s hits $m_2=1$\,kg at rest ---
find both exit speeds by the flip rule and verify $K$ survives. (AP-C)
\end{exercise}
\begin{exercise}[medium]
Conveyor heating: a $2$\,kg package accelerates to a $3$\,m/s belt over a
$1$\,m friction patch. Show the belt interface dissipates exactly the
package's gained $K$, then compute both. (F=ma)
\end{exercise}
\begin{exercise}[medium]
A $1500$\,kg car needs $25$\,kW to hold $30$\,m/s on a flat road.
Estimate the combined drag + rolling resistance at that speed, then the
power to hold $30$\,m/s up a $5\%$ grade. (engineering)
\end{exercise}
\begin{exercise}[hard]
A particle moves in $U(x)=ax^{4}-bx^{2}$ ($a,b>0$). Sketch $U$, find the
equilibria and their stability, and the energy ranges where motion is
trapped in one well vs.\ free between wells. (USAPhO-style)
\end{exercise}
\begin{exercise}[hard]
Elastic collision, unequal masses: derive the general solution
$v_1'=\frac{m_1-m_2}{m_1+m_2}v_1$,
$v_2'=\frac{2m_1}{m_1+m_2}v_1$, and verify the limits
$m_2\to\infty$ (wall) and $m_1=m_2$ (swap). (USAPhO-style)
\end{exercise}
\begin{exercise}[medium]
A spring ($k=400$\,N/m) launches a $0.5$\,kg block up a $30^{\circ}$
incline with $\mu_k=0.2$: find how far it climbs before stopping if
the spring starts compressed $10$\,cm. (AP-C FR-style)
\end{exercise}
